\subsection{Mass accretion history}

We model the halo mass accretion history (MAH) with the two-parameter form introduced by \citet{McBride2009MNRAS}. In the default normalization used for the Monte Carlo sampling,
\begin{equation}
    M_{\rm h}(z)=M_0(1+z)^\beta e^{-\gamma z},
\end{equation}
where $M_0$ is the halo mass at $z=0$. For the publication-style MAH tracks, it is more convenient to use the equivalent anchored form
\begin{equation}
    M_{\rm h}(z\mid z_{\rm ref},M_{\rm ref})=
    M_{\rm ref}
    \left(\frac{1+z}{1+z_{\rm ref}}\right)^\beta
    e^{-\gamma (z-z_{\rm ref})},
\end{equation}
which enforces $M_{\rm h}(z_{\rm ref})=M_{\rm ref}$. In the figure shown below we adopt $z_{\rm ref}=6$ and $M_{\rm ref}=10^{12}\,M_\odot$.

The corresponding accretion rate is
\begin{equation}
    \dot{M}_{\rm h}=\frac{dM_{\rm h}}{dt}
    =\frac{dM_{\rm h}}{dz}\frac{dz}{dt},
\end{equation}
with
\begin{equation}
    \frac{dM_{\rm h}}{dz}
    =
    M_{\rm h}(z)\left[\frac{\beta}{1+z}-\gamma\right].
\end{equation}
For a flat $\Lambda$CDM cosmology,
\begin{equation}
    \frac{dz}{dt}=-H(z)(1+z),
\end{equation}
\begin{equation}
    H(z)=H_0\sqrt{\Omega_m(1+z)^3+\Omega_\Lambda},
\end{equation}
so that
\begin{equation}
    \dot{M}_{\rm h}(z,\beta,\gamma)
    =
    -H(z)(1+z)M_{\rm h}(z)\left[\frac{\beta}{1+z}-\gamma\right],
\end{equation}
and the specific accretion rate is
\begin{equation}
    \frac{\dot{M}_{\rm h}}{M_{\rm h}}
    =
    H(z)\left[\gamma(1+z)-\beta\right].
\end{equation}
Throughout this work we adopt the Planck18 cosmology, with $H_0=67.4~{\rm km~s^{-1}~Mpc^{-1}}$ (equivalently $0.0689~{\rm Gyr^{-1}}$), $\Omega_m=0.315$, and $\Omega_\Lambda=0.685$.

For the Monte Carlo realizations of $(\beta,\gamma)$, our fiducial sampler follows the Appendix A prescription of \citet{McBride2009MNRAS}. In the implementation used here, this distribution is treated as a mixture of a power-law branch with $\gamma=0$ and weight $0.0466$, and a two-dimensional joint component sampled in $(\beta,y')$ and mapped to $\gamma$ using the Appendix A transformation. Both components are generated with rejection sampling. For comparison, we also construct a bivariate Gaussian approximation whose mean vector and covariance matrix are estimated from a pilot sample drawn from the original Appendix A distribution.

Each accepted $(\beta,\gamma)$ pair defines one MAH trajectory, from which we compute $M_{\rm h}(z)$, $\dot{M}_{\rm h}(z)$, and $\dot{M}_{\rm h}/M_{\rm h}$ at the target redshift. We discard non-finite realizations, cases with $M_{\rm h}(z)\le 0$, samples with $\gamma<0$, and extreme draws outside the range $-10.5<\beta<12.5$. For the anchored MAH figure, we further terminate each track once it falls below the minimum virial mass $M_{\rm min}(z)$ returned by \texttt{massfunc.M\_vir} for $\mu=0.61$ and $T_{\rm vir}=10^4~{\rm K}$. We summarize the resulting distributions using their medians and percentile ranges.

\begin{figure}[t]
    \centering
    \includegraphics[width=\linewidth]{outputs/mah_z6_Mh1e12_to_z50_mmin.pdf}
    \caption{Monte Carlo realizations of anchored halo MAHs drawn from the McBride Appendix A distribution. All tracks are normalized to $M_{\rm h}(z=6)=10^{12}\,M_\odot$. The gray curves show individual MAH realizations, while the blue dashed curve marks the minimum virial mass $M_{\rm min}(z)$. Each trajectory is truncated once it drops below this cooling threshold.}
    \label{fig:mah_z6_anchor_mmin}
\end{figure}
